\chapter{Berechnungstheoretische Analyse bedingter Unabhängigkeit in ausgewählten Graphklassen}
\label{ch:graphclasses}

Die bedingte Unabhängigkeit \(X\perp_\sigma Y\mid Z\) wird über die Kombinierbarkeit 
von \(\sigma\)-Labellings definiert. Strukturelle Eigenschaften eines AFs 
beeinflussen daher unmittelbar, welche Labellings existieren und ob zwischen 
ihnen die für eine nichttriviale Abhängigkeit erforderliche Variation auftreten kann. 
Eine besondere Rolle spielen dabei gerichtete Zyklen: Ihr Auftreten und ihre Parität 
wirken sich unterschiedlich aber oftmals grundlegend auf die betrachteten Semantiken aus. 
Dieses Kapitel untersucht deshalb, wie sich die bedingte Unabhängigkeit in bestimmten 
Graphklassen, die aufgrund ihrer vorteilhaften Eigenschaften und geringerer Komplexität 
in der Literatur betrachtet werden \cite{dung:1995:af,DUNNE2002187,handbook__of_formal_argumentation__vol__3:2024}, verhält.

Zunächst werden in Abschnitt~\ref{sec:acyc-afs} AFs ohne gerichteten Zyklus betrachtet. 
Da ihre Struktur die semantische Bewertung für
\(\sigma\in\{\mathsf{co},\mathsf{st},\mathsf{pr}\}\) 
eindeutig festlegt, gilt \(X\perp_\sigma Y\mid Z\) unabhängig von der Wahl 
der paarweise disjunkten Mengen \(X,Y,Z\). Abschnitt~\ref{sec:even-cycle-free-afs} 
schwächt diese Voraussetzung ab und erlaubt gerichtete Zyklen ungerader Länge. 
Auch in dieser Graphklasse fehlt jedoch die semantische Variation, 
die für eine nichttriviale Abhängigkeit erforderlich wäre.

In Abschnitt~\ref{sec:odd-cycle-free-afs} ändert sich die Ausgangslage grundlegend, 
da gerade Zyklen und damit mehrere Labellings zulässig bleiben. 
Der eigenständige Beitrag dieses Kapitels besteht insbesondere in der Bestimmung der 
Komplexität der \emph{stable}-, \emph{preferred}- und \emph{complete}-Unabhängigkeit 
auf dieser Graphklasse. Für die \emph{stable}-Semantik wird gezeigt, dass die bekannte 
Härtereduktion von \textcite{bluemel:2025:independence-aba} ausschließlich 
ungeradezykelfreie AFs erzeugt; das Ergebnis für die \emph{preferred}-Semantik folgt 
aus der Übereinstimmung beider Semantiken auf dieser Graphklasse. Für die \emph{complete}-Semantik 
wird eine neue Konstruktion entwickelt, welche die zusätzlich möglichen 
partiellen Labellings kontrolliert. 
Damit wird für alle drei Semantiken \(\Pi_2^{\mathrm{P}}\)-Vollständigkeit nachgewiesen.

\section{Azyklische AFs}
\label{sec:acyc-afs}
Ein AF \(F=(A,R)\) heißt azyklisch, wenn der gerichtete Graph \((A,R)\) keinen gerichteten Zyklus enthält. 
Azyklische AFs bilden eine besonders restriktive Graphklasse, für sie fallen alle betrachteten 
Semantiken auf das eindeutige \emph{grounded}-Labelling zusammen \cite{dung:1995:af,caminada:2009:logical-account}.
Im Folgenden wird gezeigt, dass diese strukturelle Einschränkung 
eine unmittelbare Aussage über \(X\perp_\sigma Y\mid Z\) erlaubt.
\begin{Proposition}
\label{prop:acyc-af-ci-st-top}
    Sei \(F=(A,R)\) ein azyklisches AF und \(\sigma\) eine Semantik mit
    \(\sigma \in \{\mathsf{co},\mathsf{st},\mathsf{pr}\}\). Dann gilt für alle paarweise disjunkten
    \(X,Y,Z \subseteq A\), dass \(X\perp_\sigma Y\mid Z\).
\end{Proposition}

\begin{proof}
\label{pr:acyc-af-ci-st-top-true}
    Sei \(F=(A,R)\) ein azyklisches AF. Für azyklische AFs fallen die \emph{complete}-, 
    \emph{preferred}- und \emph{stable}-Semantik mit der \emph{grounded}-Semantik zusammen. 
    Da das \emph{grounded}-Labelling eindeutig ist, existiert genau ein Labelling \(\lambda\),
     sodass \(\Lambda_{\mathrm{co}}(F)=\Lambda_{\mathrm{pr}}(F)=\Lambda_{\mathrm{st}}(F)=\{\lambda\}\).
    Sei nun \(\sigma \in \{\mathsf{co},\mathsf{st},\mathsf{pr}\}\). Dann gilt insbesondere \(|\Lambda_{\sigma}(F)|=1\).
     Nach Lemma~\ref{lem:independence-trivial-labelling-set} folgt für alle 
     paarweise disjunkten \(X,Y,Z \subseteq A\), dass \(X \perp_\sigma Y \mid Z\) gilt.
\end{proof}
\section{Geradezykelfreie AFs}
\label{sec:even-cycle-free-afs}
Geradezykelfreie AFs schließen gerichtete Zyklen gerader Länge aus, erlauben jedoch ungerade Zyklen. 
Sie bilden daher eine echt umfassendere Graphklasse als die azyklischen AFs. 
Dennoch begrenzt das Fehlen gerader Zyklen die Menge der möglichen Labellings so stark \cite{dvorak:2012:computational-aspects}, 
dass sich erneut allgemeine Aussagen über \(X\perp_\sigma Y\mid Z\) ableiten lassen.
\begin{Proposition}
\label{prop:even-cycle-free-af-ci-st-top}
    Sei \(F=(A,R)\) ein AF ohne gerade Zyklen und \(\sigma\) eine Semantik mit
    \(\sigma \in \{\mathsf{co},\mathsf{st},\mathsf{pr}\}\). Dann gilt für alle paarweise disjunkten
    \(X,Y,Z \subseteq A\), dass \(X\perp_\sigma Y\mid Z\).
\end{Proposition}

\begin{proof}
\label{even-cycle-free-af-ci-top-true}
    Sei \(F=(A,R)\) ein AF ohne gerade Zyklen. Für diese Klasse fallen die \emph{complete}- und die \emph{preferred}-Labellings mit dem \emph{grounded}-Labelling zusammen. 
    Da das \emph{grounded}-Labelling eindeutig ist, gilt \(|\Lambda_{\mathrm{co}}(F)|= |\Lambda_{\mathrm{pr}}(F)|=1\).
    Für die \emph{stable}-Semantik ist das \emph{grounded}-Labelling der einzige mögliche Kandidat. 
    Daher gilt entweder \(\Lambda_{\mathrm{st}}(F)=\emptyset\), falls das \emph{grounded}-Labelling nicht \emph{stable} ist, oder \(\Lambda_{\mathrm{st}}(F)=\{\lambda\}\), 
    falls das \emph{grounded}-Labelling \emph{stable} ist. In beiden Fällen gilt \(|\Lambda_{\mathrm{st}}(F)| \leq 1\).
    Somit gilt für jedes \(\sigma \in \{\mathsf{co},\mathsf{st},\mathsf{pr}\}\), dass \(|\Lambda_{\sigma}(F)| \leq 1\). 
    Nach Lemma~\ref{lem:independence-trivial-labelling-set} folgt für alle paarweise disjunkten \(X,Y,Z \subseteq A\), dass \(X \perp_\sigma Y \mid Z\) gilt.
\end{proof}

Auch in AFs ohne Zyklen gerader Länge beruht die Allgemeingültigkeit der 
bedingten Unabhängigkeit damit lediglich auf der fehlenden Variation
der Labellings.  
\section{Ungeradezykelfreie AFs}
\label{sec:odd-cycle-free-afs}
In den bisher betrachteten Graphklassen entstand bedingte Unabhängigkeit aus dem Fehlen 
alternativer Labellings. Mit den ungeradezykelfreien AFs ändert sich diese Ausgangslage: 
Zyklen bleiben nun zulässig, sofern sie eine gerade Länge besitzen. Damit bleiben gerade 
jene Rückkopplungen erhalten, die größere Labellingvariationen ermöglichen. 

Dennoch fallen in ungeradezykelfreien AFs die \emph{stable}- und die \emph{preferred}-Semantik zusammen \cite{DUNNE2002187}.
Die folgenden Ergebnisse zeigen, dass die strukturelle Einschränkung keine weitere Verringerung 
der Komplexität bewirkt. Während sich dies für die \emph{stable}- und die \emph{preferred}-Semantik 
aus der bekannten Reduktion von \textcite{bluemel:2025:independence-aba} und ihrer Graphstruktur ableiten lässt, 
erfordert die \emph{complete}-Semantik eine neue Konstruktion.

Die Beschränkung auf ungeradezykelfreie AFs senkt die Komplexität der \emph{stable}-Unabhängigkeit nicht.

\begin{Proposition}
    \label{prop:ocf-af-ci-st-pi^2}
    Die \emph{stable}-Unabhängigkeit ist auch auf ungeradezykelfreien AFs
    \(\Pi_2^{\mathrm{P}}\)-vollständig.
\end{Proposition}

\begin{proof}
    \label{pr:ocf-af-ci-st-pi^2-is-correct}
    Die Zugehörigkeit zu \(\Pi_2^{\mathrm{P}}\) folgt aus dem entsprechenden
    Ergebnis für allgemeine AFs
    \parencite[Proposition~4.10]{bluemel:2025:independence-aba}.
    Für die Härte betrachten wir die direkte AF-Konstruktion aus
    \textcite[Proposition~C.4]{bluemel:2025:independence-aba}.
    Für jedes \(v\in X\cup Y\cup\{x_0\}\) bildet
    \(\{v,\overline v\}\) eine stark zusammenhängende Komponente.
    Daneben bilden \(\{\varphi,t\}\) sowie jedes Singleton
    \(\{c_k\}\) die übrigen stark zusammenhängenden Komponenten.
    Jeder Angriff zwischen verschiedenen Komponenten verläuft gemäß der
    Richtung:
    \(
        \{v,\overline v\}
        \rightarrow
        \{c_k\}
        \rightarrow
        \{\varphi,t\}.
    \)
    Ein gerichteter Zyklus kann daher keine Komponente verlassen. Innerhalb
    der nichttrivialen Komponenten existieren ausschließlich die Zweierzyklen
    \(v\leftrightarrow\overline v\) und
    \(\varphi\leftrightarrow t\).Außerdem besitzen die Klauselargumente keine
    Selbstangriffe. Somit hat jeder gerichtete Zyklus gerade Länge und jedes
    durch die Reduktion erzeugte AF ist ungeradezykelfrei. Die
    \(\Pi_2^{\mathrm{P}}\)-Härte gilt folglich auch auf dieser Graphklasse.
\end{proof}

\begin{figure}
    \centering
    \resizebox{\textwidth}{!}{
        \tikzset{
    arg/.style={
        draw,
        circle,
        fill=gray!15,
        inner sep=1pt,
        minimum size=.5cm
    }
}

\begin{tikzpicture}[>=stealth,xscale=1, yscale=0.7]
  \small

  % Zielargumente
  \node[arg] (phi) at (0,0.5) {$\varphi$};
  \node[arg] (t)   at (1.8,0.5) {$t$};

  % Klauselargumente
  \node[arg] (cl1) at (-3.0,-1.3) {$cl_1$};
  \node[arg] (cl2) at ( 0.0,-1.3) {$cl_2$};
  \node[arg] (cl3) at ( 3.0,-1.3) {$cl_3$};

  % Literalargumente
  \node[arg] (x0)  at (-6.3,-4.0) {$x_0$};
  \node[arg] (nx0) at (-4.9,-4.0) {$\bar{x}_0$};
  \node[arg] (x1)  at (-3.2,-4.0) {$x_1$};
  \node[arg] (nx1) at (-1.8,-4.0) {$\bar{x}_1$};
  \node[arg] (x2)  at ( 0.0,-4.0) {$x_2$};
  \node[arg] (nx2) at ( 1.4,-4.0) {$\bar{x}_2$};
  \node[arg] (y1)  at ( 3.2,-4.0) {$y_1$};
  \node[arg] (ny1) at ( 4.6,-4.0) {$\bar{y}_1$};

  % Gegenseitige Angriffe komplementärer Literale
  \path[attack, <->, thick]
    (x0) edge (nx0)
    (x1) edge (nx1)
    (x2) edge (nx2)
    (y1) edge (ny1)
    (phi) edge (t);

  % Erfüllte Literale greifen die zugehörigen Klauseln an
  \path[attack, thick]
    (x0)  edge (cl1)
          edge (cl2)
          edge (cl3)
    (x2)  edge (cl1)
    (y1)  edge (cl1)
    (nx2) edge (cl2)
          edge (cl3)
    (ny1) edge (cl2)
    (nx1) edge (cl2)
    (x1)  edge (cl3);

  % Klauseln greifen das Formelargument an
  \path[attack, thick]
    (cl1) edge (phi)
    (cl2) edge (phi)
    (cl3) edge (phi);
\end{tikzpicture}
    }
    \caption[Reduktion für \emph{stable}-Unabhängigkeit]{
    AF-Repräsentation einer Instanz der Reduktion aus Proposition~C.4
    für die Formel \(\Psi=\forall X\,\exists Y\,\varphi\).
    Die Matrix \(\varphi\) ist durch die Klauselmenge
    \(\{\{x_0,x_2,y_1\},
    \{x_0,\overline{x_2},\overline{y_1},\overline{x_1}\},
    \{x_0,\overline{x_2},x_1\}\}\)
    gegeben. Darstellung adaptiert nach
    \cite{bluemel:2025:independence-aba}.
    }   
    \label{fig:reduktion-stable}
\end{figure}

Abbildung~\ref{fig:reduktion-stable} zeigt eine konkrete Instanz der Reduktion.
Das vorangegangene Ergebnis lässt sich für ungeradezykelfreien AFs unmittelbar auf die \emph{preferred}-Unabhängigkeit übertragen. 
Maßgeblich ist die Gleichheit \(\Lambda_{\mathrm{st}}(F)=\Lambda_{\mathrm{pr}}(F)\), 
aufgrund derer beide Unabhängigkeitsbegriffe auf derselben Menge von Labellings ausgewertet werden.
\begin{Proposition}
    \label{prop:ocf-af-ci-pr-pi^2-c}
    Die \emph{preferred}-Unabhängigkeit ist auch in ungeradezykelfreien AFs \(\Pi_2^{\mathrm{P}}\)-vollständig.
\end{Proposition}

\begin{proof}
    \label{pr:ocf-af-ci-pr-pi^2-c-is-correct}
    Dies folgt direkt aus Proposition~\ref{prop:ocf-af-ci-st-pi^2} und der Tatsache, dass in ungeradezykelfreien AFs \(\Lambda_{\mathrm{st}}(F)=\Lambda_{\mathrm{pr}}(F)\).
\end{proof}

Für die \emph{complete}-Unabhängigkeit ist ein entsprechender Übertrag nicht möglich, 
da \(\Lambda_{\mathrm{co}}(F)\) zusätzliche, nichtmaximale Labellings enthalten kann. 
Insbesondere können Variablenpaare mit \(\mathsf{und}\) gelabelt sein, weshalb eine 
modifizierte Konstruktion benötigt wird, die totale und partielle Belegungen unterscheidet.
\begin{Proposition}
    \label{prop:ocf-af-ci-co-pi^2-c}
    Die \emph{complete}-Unabhängigkeit ist auch in ungeradezykelfreien AFs \(\Pi_2^{\mathrm{P}}\)-vollständig.
\end{Proposition}

\begin{proof}
    \label{pr:ocf-af-ci-co-pi^2-c-is-correct}
    Wir modifizieren die Konstruktion von \textcite{bluemel:2025:independence-aba}.
    Reduziert wird vom \(\Pi_2^{\mathrm{P}}\)-vollständigen Problem
    \(\mathrm{QBF}^{2}_{\forall}\). Eine Instanz dieses Problems ist
    eine quantifizierte boolesche Formel der Form
    \(\Psi=\forall X\exists Y\,\varphi(X,Y), \)
    wobei \(\varphi=\bigwedge_{k=1}^{m}C_k\) in konjunktiver
    Normalform vorliegt. Gefragt ist, ob für jede Belegung der
    Variablen aus \(X\) eine Belegung der Variablen aus \(Y\)
    existiert, welche \(\varphi\) erfüllt.

    \paragraph{Konstruktion.}
    Wir führen eine neue universelle
    Variable \(x_0\) ein und setzen
    \( X'=X\cup\{x_0\}~
        \text{und}~
        \widehat{\varphi}
        =
        \bigwedge_{k=1}^{m}(x_0\lor C_k).\)
    Die Formel \(\forall X\exists Y\,\varphi\) ist genau dann gültig,
    wenn die Formel \(\forall X'\exists Y\,\widehat{\varphi}\) gültig ist:
    Für \(x_0=\top\) ist jede Klausel von \(\widehat{\varphi}\) erfüllt,
    während \(\widehat{\varphi}\) für \(x_0=\bot\) mit \(\varphi\)
    übereinstimmt.
    Aus \(\Psi\) konstruieren wir das AF
    \(F_{\Psi}=(A_{\Psi},R_{\Psi})\) mit
    \[
    \begin{split}
        A_{\Psi}={}&
        \{x_i,\overline{x_i},d_i\mid x_i\in X'\}
        \cup
        \{y_j,\overline{y_j}\mid y_j\in Y\}\\
        &{}\cup
        \{c_1,\ldots,c_m,\varphi,t,s,g\}.
    \end{split}
    \]
    Die Angriffsmenge ist wie folgt definiert:
    \[
    \begin{split}
        R_{\Psi}={}&
        \{(x_i,\overline{x_i}),(\overline{x_i},x_i)
            \mid x_i\in X'\}\\
        &{}\cup
        \{(y_j,\overline{y_j}),(\overline{y_j},y_j)
            \mid y_j\in Y\}\\
        &{}\cup
        \{(\ell,c_k)
            \mid \ell\text{ kommt in der Klausel }
            x_0\lor C_k\text{ vor}\}\\
        &{}\cup
        \{(c_k,\varphi)\mid 1\leq k\leq m\}\\
        &{}\cup
        \{(\varphi,t),(t,\varphi)\}\\
        &{}\cup
        \{(x_i,d_i),(\overline{x_i},d_i),(d_i,s)
            \mid x_i\in X'\}\\
        &{}\cup
        \{(s,g),(g,\varphi)\}\\
        &{}\cup
        \{(g,c_k)\mid 1\leq k\leq m\}.
    \end{split}
    \]
    Die Konstruktion ist damit polynomiell in der
    Größe von \(\Psi\).

    Abbildung~\ref{fig:reduktion-complete} veranschaulicht die Konstruktion
    und zeigt insbesondere, wie die Kontrollstruktur aus \(d_i\), \(s\) und
    \(g\) totale von partiellen Belegungen trennt. Die Konstruktion
    berücksichtigt damit die zusätzlichen \emph{complete}-Labellings, während
    alle gerichteten Zyklen weiterhin auf die Zweierzyklen innerhalb der
    Variablenpaare und der Komponente \(\{\varphi,t\}\) beschränkt bleiben.

    Als Instanz der \emph{complete}-Unabhängigkeit geben wir
    \(
        \{\varphi\}
        \perp_{\mathrm{co}}
        X'
        \mid
        \{s\}
    \)
    aus.

    \paragraph{Vorbeobachtungen.}
    Sei \(\lambda\in\Lambda_{\mathrm{co}}(F_{\Psi})\) ein beliebiges
    \emph{complete}-Labelling. Für jedes \(x_i\in X'\) gilt
    \[
        \bigl(\lambda(x_i),\lambda(\overline{x_i})\bigr)
        \in
        \{(\mathsf{in},\mathsf{out}),
          (\mathsf{out},\mathsf{in}),
          (\mathsf{und},\mathsf{und})\}.
    \]
    In den ersten beiden Fällen ist die zugehörige Variable total
    belegt und im dritten Fall ist sie unentschieden belegt. Da \(d_i\) von beiden
    Literalargumenten angegriffen wird, gilt \(\lambda(d_i)=\mathsf{out}\) genau
    im totalen Fall und \(\lambda(d_i)=\mathsf{und}\) im unentschiedenen Fall.
    Insbesondere erhält kein \(d_i\) das Label \(\mathsf{in}\). Folglich gilt
    \(\lambda(s)=\mathsf{in}\) genau dann, wenn alle Variablenpaare für \(X'\)
    total gelabelt sind; andernfalls ist \(\lambda(s)=\mathsf{und}\). Da \(s\)
    der einzige Angreifer von \(g\) ist, gilt entsprechend
    \(\lambda(g)=\mathsf{out}\) beziehungsweise \(\lambda(g)=\mathsf{und}\). Im
    partiellen Fall gilt also \(\lambda(g)=\mathsf{und}\). Da \(g\) jedes
    \(c_k\) und \(\varphi\) angreift, kann keines dieser Argumente mit
    \(\mathsf{in}\) gelabelt sein. Auf diese Weise wird die eigentliche
    Formelauswertung ausgesetzt.

    \paragraph{Zu zeigende Aussage.}
    Wir zeigen die folgende Äquivalenz:
    \begin{equation}
        \Psi\text{ ist genau dann gültig, wenn }
        \{\varphi\}\perp_{\mathrm{co}}X'\mid\{s\}
        \text{ in }F_{\Psi}\text{ gilt.}
        \tag{$*$}
    \end{equation}

    \paragraph{Beweis von \((*)\).}
    Sei zunächst \(\Psi\) gültig und seien
    \(\lambda_1,\lambda_2\in\Lambda_{\mathrm{co}}(F_{\Psi})\) mit
    \(\lambda_1(s)=\lambda_2(s)\). Wir konstruieren ein
    \emph{complete}-Labelling \(\lambda_3\), das das Label von \(\varphi\) aus
    \(\lambda_1\) und das Labelling von \(X'\) aus \(\lambda_2\)
    übernimmt.
    Falls \(\lambda_1(s)=\lambda_2(s)=\mathsf{und}\), folgt aus den
    Vorbeobachtungen, dass \(\lambda_1(g)=\lambda_2(g)=\mathsf{und}\) gilt.
    Da \(g\) das Argument \(\varphi\) angreift, kann \(\varphi\) in
    \(\lambda_1\) nicht mit \(\mathsf{in}\) gelabelt sein: Dafür müssten alle
    Angreifer von \(\varphi\) mit \(\mathsf{out}\) gelabelt sein. Für jedes
    \(x_i\in X'\) übernehmen wir die Labels des Variablenpaars
    \(x_i,\overline{x_i}\) aus \(\lambda_2\) und jedes Variablenpaar
    \(y_j,\overline{y_j}\) labeln wir mit \((\mathsf{und},\mathsf{und})\). Die dadurch
    bestimmten Labels der Argumente \(d_i\) übernehmen wir ebenfalls und 
    setzen ferner \(\lambda_3(s)=\lambda_3(g)=\mathsf{und}\). Jedes
    Klauselargument \(c_k\), das von einem mit \(\mathsf{in}\) gelabelten
    Literal angegriffen wird, erhält in \(\lambda_3\) das Label \(\mathsf{out}\);
    alle übrigen Klauselargumente erhalten das Label \(\mathsf{und}\). Wegen
    \(\lambda_3(g)=\mathsf{und}\) sind genau diese beiden Fälle möglich.
    Abschließend wählen wir
    \(\lambda_3(\varphi)=\mathsf{out}\) und \(\lambda_3(t)=\mathsf{in}\), falls
    \(\lambda_1(\varphi)=\mathsf{out}\) gilt. Falls
    \(\lambda_1(\varphi)=\mathsf{und}\) gilt, wählen wir dagegen
    \(\lambda_3(\varphi)=\lambda_3(t)=\mathsf{und}\). Damit übernimmt
    \(\lambda_3\) in beiden Fällen das Label von \(\varphi\) aus
    \(\lambda_1\).
    Falls \(\lambda_1(s)=\lambda_2(s)=\mathsf{in}\), kodiert
    \(\lambda_2|_{X'}\) eine totale Belegung \(\alpha'\) von
    \(X'\). Wegen der Gültigkeit von
    \(\forall X'\exists Y\,\widehat{\varphi}\) existiert eine
    Belegung \(\beta\) von \(Y\) mit
    \(\widehat{\varphi}(\alpha',\beta)=\top\). Für jedes \(x_i\in X'\)
    übernehmen wir die Labels des Variablenpaars
    \(x_i,\overline{x_i}\) aus \(\lambda_2\) und labeln die
    \(Y\)-Argumente gemäß \(\beta\). Damit sind alle \(d_i\) mit
    \(\mathsf{out}\), das Argument \(s\) mit \(\mathsf{in}\) und das Argument \(g\) mit
    \(\mathsf{out}\) gelabelt. Außerdem wird jedes
    Klauselargument von einem \(\mathsf{in}\)-gelabelten Literal angegriffen
    und somit mit \(\mathsf{out}\) gelabelt. Daher können die drei möglichen
    Labels von \(\varphi\) durch
    \(
        \bigl(\lambda_3(\varphi),\lambda_3(t)\bigr)\in
        \{(\mathsf{in},\mathsf{out}),
          (\mathsf{out},\mathsf{in}),
          (\mathsf{und},\mathsf{und})\}
    \)
    realisiert werden. Wir wählen das Paar, für das
    \(\lambda_3(\varphi)=\lambda_1(\varphi)\) gilt. Somit existiert in
    beiden Fällen das
    erforderliche Labelling \(\lambda_3\), und es gilt
    \(\{\varphi\}\perp_{\mathrm{co}}X'\mid\{s\}\).

    Für die umgekehrte Richtung zeigen wir die Kontraposition. Sei nun
    \(\Psi\) nicht gültig. Dann existiert eine Belegung
    \(\alpha\) der Variablen aus \(X\), sodass
    \(\varphi(\alpha,\beta)=\bot\) für jede Belegung \(\beta\) der
    Variablen aus \(Y\) gilt.
    Wir erweitern \(\alpha\) durch \(x_0=\bot\) zu einer totalen
    Belegung \(\alpha'\) von \(X'\) und wählen eine beliebige totale
    Belegung der Variablen aus \(Y\). Für das zu konstruierende
    Labelling \(\lambda_2\) labeln wir die Variablenpaare gemäß diesen
    Belegungen. Dadurch sind alle \(d_i\) mit \(\mathsf{out}\), das Argument
    \(s\) mit \(\mathsf{in}\) und das Argument \(g\) mit \(\mathsf{out}\) gelabelt. Ein
    Klauselargument erhält das Label \(\mathsf{out}\), wenn die zugehörige
    Klausel erfüllt ist, und andernfalls das Label \(\mathsf{in}\). Da
    \(x_0=\bot\) gilt und keine Belegung von \(Y\) die Formel
    \(\varphi\) unter \(\alpha\) erfüllt, ist mindestens ein
    Klauselargument mit \(\mathsf{in}\) gelabelt. Daher gilt notwendigerweise
    \(\lambda_2(\varphi)=\mathsf{out}\) und folglich \(\lambda_2(t)=\mathsf{in}\).
    Die so festgelegten Labels bilden ein \emph{complete}-Labelling,
    das \(\alpha'\) kodiert und für das \(\lambda_2(s)=\mathsf{in}\) gilt.
    Separat konstruieren wir ein \emph{complete}-Labelling \(\lambda_1\)
    mit \(\lambda_1(s)=\mathsf{in}\) und \(\lambda_1(\varphi)=\mathsf{in}\). Dazu wählen
    wir eine totale Belegung der Variablen aus \(X'\cup Y\) mit
    \(x_0=\top\) und labeln die Variablenpaare entsprechend. Das
    Argument \(x_0\) ist dann mit \(\mathsf{in}\) gelabelt und greift alle
    Klauselargumente an, sodass diese das Label \(\mathsf{out}\) erhalten. Weil
    alle Variablenpaare total gelabelt sind, gilt außerdem
    \(\lambda_1(s)=\mathsf{in}\) und \(\lambda_1(g)=\mathsf{out}\). Schließlich setzen
    wir \(\lambda_1(\varphi)=\mathsf{in}\) und \(\lambda_1(t)=\mathsf{out}\).

    Angenommen, es gäbe ein \emph{complete}-Labelling \(\lambda_3\),
    für das \(\lambda_3(\varphi)=\lambda_1(\varphi)=\mathsf{in}\),
    \(\lambda_3|_{X'}=\lambda_2|_{X'}\) und
    \(\lambda_3(s)=\lambda_1(s)=\lambda_2(s)=\mathsf{in}\) gilt. Aus
    \(\lambda_3(\varphi)=\mathsf{in}\) folgt, dass alle Klauselargumente \(c_k\)
    mit \(\mathsf{out}\) gelabelt sind. Wegen \(\lambda_3(s)=\mathsf{in}\) gilt auch
    \(\lambda_3(g)=\mathsf{out}\). Daher muss
    jedes \(c_k\) von einem mit \(\mathsf{in}\) gelabelten Literalargument
    angegriffen werden. Die mit \(\mathsf{in}\) gelabelten \(Y\)-Literale
    bestimmen eine partielle Belegung von \(Y\), die sich beliebig zu
    einer totalen Belegung \(\beta\) erweitern lässt. Jede solche
    Erweiterung erfüllt zusammen mit \(\alpha'\) alle Klauseln von
    \(\widehat{\varphi}\). Wegen \(x_0=\bot\) folgt daraus
    \(\varphi(\alpha,\beta)=\top\), im Widerspruch zur Wahl von
    \(\alpha\). Ein solches \(\lambda_3\) existiert daher nicht, und
    \(\{\varphi\}\not\perp_{\mathrm{co}}X'\mid\{s\}\) gilt.

    Damit ist die Aussage \((*)\) bewiesen.

    \paragraph{Graphklassenzugehörigkeit.}
    Das konstruierte AF \(F_{\Psi}\) ist ungeradezykelfrei. 
    Seine einzigen nichttrivialen stark zusammenhängenden Komponenten sind
    \(\{x_i,\overline{x_i}\}\) für \(x_i\in X'\),
    \(\{y_j,\overline{y_j}\}\) für \(y_j\in Y\) sowie
    \(\{\varphi,t\}\).
    Die Argumente \(d_i\), \(s\), \(g\) und \(c_k\) bilden jeweils
    einelementige stark zusammenhängende Komponenten.
    Zwischen diesen Komponenten verlaufen die Angriffe ausschließlich in den
    Richtungen
    \(\{x_i,\overline{x_i}\}\to d_i\to s\to g\),
    \(\{x_i,\overline{x_i}\},\{y_j,\overline{y_j}\}\to c_k\),
    \(g\to c_k\) und
    \(g,c_k\to\{\varphi,t\}\).
    Folglich liegt jeder gerichtete Zyklus vollständig in einer der
    nichttrivialen Komponenten. Diese enthalten jeweils ausschließlich den
    Zweierzyklus
    \(x_i\leftrightarrow\overline{x_i}\),
    \(y_j\leftrightarrow\overline{y_j}\) beziehungsweise
    \(\varphi\leftrightarrow t\).
    Somit besitzt jeder gerichtete Zyklus gerade Länge und
    \(F_{\Psi}\) ist ungeradezykelfrei.
\end{proof}

\begin{figure}
    \centering
    \resizebox{\textwidth}{!}{
        \tikzset{
    arg/.style={
        draw,
        circle,
        fill=gray!15,
        inner sep=1pt,
        minimum size=.5cm
    },
    auxiliary/.style={
        arg,
        fill=gray!30
    },
    attack/.style={
        ->,
        thick,
        >=stealth
    }
}

\begin{tikzpicture}[>=stealth,xscale=1, yscale=0.7]
\small

    % Formel-, Klausel- und Hilfsargumente
    \path
        node[arg]       at ( 1.2,  .3) (phi) {$\varphi$}
        node[arg]       at ( 3.2,  .3) (t)   {$t$}
        node[arg]       at (-1.1,-2.0) (cl1)  {$cl_1$}
        node[arg]       at ( 1.8,-2.0) (cl2)  {$cl_2$}
        node[auxiliary] at ( 4.8,-2.8) (g)   {$g$};

    % Literalargumente
    \path
        node[arg] at (-4.9,-5.0) (x0)  {$x_0$}
        node[arg] at (-3.5,-5.0) (nx0) {$\bar{x}_0$}
        node[arg] at (-1.9,-5.0) (x1)  {$x_1$}
        node[arg] at (-0.5,-5.0) (nx1) {$\bar{x}_1$}
        node[arg] at ( 1.5,-5.0) (y1)  {$y_1$}
        node[arg] at ( 2.9,-5.0) (ny1) {$\bar{y}_1$};

    % Unteres Hilfsgadget
    \path
        node[auxiliary] at (-4.2,-7.3) (d0) {$d_0$}
        node[auxiliary] at (-1.2,-7.3) (d1) {$d_1$}
        node[auxiliary] at (-2.7,-9.3) (s)  {$s$};

    % Symmetrische Angriffe werden durch eine Kante mit zwei
    % Pfeilspitzen dargestellt.
    \path[attack,<->]
        (x0)  edge (nx0)
        (x1)  edge (nx1)
        (y1)  edge (ny1)
        (phi) edge (t);

    % Literale greifen die Klauseln an, in denen sie vorkommen
    \path[attack]
        (x0)  edge (cl1)
              edge (cl2)
        (x1)  edge (cl1)
        (y1)  edge (cl1)
        (nx1) edge (cl2)
        (ny1) edge (cl2);

    % Klauselargumente greifen das Formelargument an
    \path[attack]
        (cl1) edge (phi)
        (cl2) edge (phi);

    % d- und s-Gadget
    \path[attack]
        (x0)  edge (d0)
        (nx0) edge (d0)
        (x1)  edge (d1)
        (nx1) edge (d1)
        (d0)  edge (s)
        (d1)  edge (s)
        (s)   edge[out=-5,in=-90,looseness=1] (g);

    % Angriffe des g-Arguments
    \path[attack]
        (g) edge[out=150,in=30,looseness=1.15] (cl1)
            edge[out=165,in=0] (cl2)
            edge[out=105,in=-20,looseness=1.05] (phi);
\end{tikzpicture}

    }
    \caption[Reduktion für \emph{complete}-Unabhängigkeit]{
    AF-Repräsentation einer Instanz der Reduktion aus
    Proposition~\ref{prop:ocf-af-ci-co-pi^2-c}
    für die Formel
    \(\Psi=\forall x_1\,\exists y_1\,\varphi\) mit
    \(\varphi=(x_1\lor y_1)
    \land(\overline{x_1}\lor\overline{y_1})\).
    Die daraus konstruierte Formel \(\widehat{\varphi}\)
    ist durch die Klauselmenge
    \(\{\{x_0,x_1,y_1\},
    \{x_0,\overline{x_1},\overline{y_1}\}\}\)
    gegeben. Die dunkler hinterlegten Knoten kennzeichnen
    die zusätzlichen Hilfsargumente \(d_0,d_1,s\) und \(g\).
    }
\label{fig:reduktion-complete}
\end{figure}



\begin{table}[H]
    \myfloatalign
    \small
    \renewcommand{\arraystretch}{1.15}
    \begin{tabular}{@{}lccc@{}}
        \toprule
        AF-Klasse
            & \(\mathsf{co}\)
            & \(\mathsf{st}\)
            & \(\mathsf{pr}\) \\
        \midrule
        azyklisch
            & trivial & trivial & trivial \\
        geradezykelfrei
            & trivial & trivial & trivial \\
        ungeradezykelfrei
            & \(\Pi_2^{\mathrm{P}}\)-c
            & \(\Pi_2^{\mathrm{P}}\)-c
            & \(\Pi_2^{\mathrm{P}}\)-c \\
        \bottomrule
    \end{tabular}
    \caption[Komplexität der bedingten Unabhängigkeit in ausgewählten Graphklassen]{
        Komplexität der bedingten \(\sigma\)-Unabhängigkeit in den
        betrachteten Graphklassen. Dabei bezeichnet \(C\)-c die
        Vollständigkeit für die Komplexitätsklasse \(C\).
    }
    \label{tab:komplexitaet-graphklassen}
\end{table}

Tabelle~\ref{tab:komplexitaet-graphklassen} zeigt: Die bedingte
\(\sigma\)-Unabhängigkeit ist auf azyklischen und geradezykelfreien AFs
für alle drei betrachteten Semantiken trivial; jede Anfrage wird positiv
beantwortet. Auf ungeradezykelfreien AFs ist das Entscheidungsproblem
dagegen für alle drei Semantiken \(\Pi_2^{\mathrm{P}}\)-vollständig und liegt
damit auf der zweiten Ebene der polynomiellen Hierarchie.





\iffalse
\section{Bipartite AFs}
\begin{Proposition}
\label{prop:bip-st-ind-hard}
    \(stable\)-Unabhängigkeit ist auch in bipartiten AFs \(\Pi_2^{\mathrm{P}}\)-vollständig.
\end{Proposition}

\begin{proof}
    Wir reduzieren vom \(\Pi_2^{\mathrm{P}}\)-vollständigen Problem
    \textsc{Balanced Monotone \(\forall\exists\)-3-SAT-\((1,1,2,2)\)}.

    \paragraph{Konstruktion.}
Sei
\(
    \Psi=\forall X\exists Y\,\varphi
\)
eine Instanz von \textsc{Balanced Monotone \(\forall\exists\)-3-SAT-\((1,1,2,2)\)}.
Die Klauselmenge von \(\varphi\) zerfällt in die positiven Klauseln
\(\mathcal{C}^{+}\) und die negativen Klauseln \(\mathcal{C}^{-}\).
Wir konstruieren das AF
\(
    F_{\Psi}=(\mathcal{A}_{\Psi},R_{\Psi}).
\)
Die Argumentmenge ist gegeben durch
\[
\begin{aligned}
    \mathcal{A}_{\Psi}
    ={}&
    \{v,\overline{v}\mid v\in X\cup Y\cup\{p,n\}\}\\
    &\cup
    \{c_C\mid C\in\mathcal{C}^{+}\cup\mathcal{C}^{-}\}\\
    &\cup
    \{\varphi^{+},\varphi^{-},t^{+},t^{-}\},
\end{aligned}
\]
wobei \(p\) und \(n\) neu eingeführte Variablen sind.
Die Angriffsrelation ist
\[
\begin{aligned}
    R_{\Psi}
    ={}&
    \{(v,\overline{v}),(\overline{v},v)
        \mid v\in X\cup Y\cup\{p,n\}\}\\
    &\cup
    \{(v,c_C)
        \mid C\in\mathcal{C}^{+},\,v\in C\}\\
    &\cup
    \{(\overline{v},c_C)
        \mid C\in\mathcal{C}^{-},\,\neg v\in C\}\\
    &\cup
    \{(p,c_C)\mid C\in\mathcal{C}^{+}\}\\
    &\cup
    \{(\overline{n},c_C)\mid C\in\mathcal{C}^{-}\}\\
    &\cup
    \{(c_C,\varphi^{+})\mid C\in\mathcal{C}^{+}\}\\
    &\cup
    \{(c_C,\varphi^{-})\mid C\in\mathcal{C}^{-}\}\\
    &\cup
    \{(\varphi^{+},t^{+}),(t^{+},\varphi^{+})\}\\
    &\cup
    \{(\varphi^{-},t^{-}),(t^{-},\varphi^{-})\}.
\end{aligned}
\]
Als \(stable\)-Unabhängigkeitsinstanz wählen wir
\(\{\varphi^{+},\varphi^{-}\} \perp_{\mathrm{st}} X \cup \{p,n\} \mid \emptyset\).

Die gegenseitigen Angriffe
\(
    v\leftrightarrow\overline{v}
\)
erzwingen unter stabiler Semantik, dass genau eines der beiden
Literalargumente \(\mathsf{in}\) ist. Jedes stabile Labelling kodiert damit eine
totale Wahrheitsbelegung.
Ein positives Klauselargument \(c_C\) wird von den positiven
Literalargumenten seiner Klausel sowie von \(p\) angegriffen. Ein negatives
Klauselargument wird von den entsprechenden negativen Literalargumenten
sowie von \(\overline{n}\) angegriffen. Wenn
\(
    p=\mathsf{out}
    \text{und}
    n=\mathsf{in},
\)
dann werden die Klauselargumente ausschließlich durch die ursprünglichen Literale
kontrolliert werden.
Falls
\(
    p=\mathsf{in}
    \text{und}
    n=\mathsf{out}
\)
sind \(p\) und \(\overline{n}\) jeweils \(\mathsf{in}\). Dadurch werden alle positiven
beziehungsweise negativen Klauselargumente angegriffen. Somit ist
\(
    \varphi^{+}=\varphi^{-}=\mathsf{in}
\)
unabhängig von der ursprünglichen Formel separat realisierbar.
Die Argumente \(t^{+}\) und \(t^{-}\) erlauben es, die Labels von
\(\varphi^{+}\) und \(\varphi^{-}\) unabhängig auf \(\mathsf{in}\) oder \(\mathsf{out}\) zu
setzen, sofern die zugehörigen Klauselargumente \(\mathsf{out}\) sind.

\paragraph{Graphklassenzugehörigkeit.}
Definiere
\[
\begin{aligned}
    L={}&
    \{v\mid v\in X\cup Y\cup\{p,n\}\}\\
    &\cup
    \{c_C\mid C\in\mathcal{C}^{-}\}
    \cup
    \{\varphi^{+},t^{-}\},
\end{aligned}
\]
und
\[
\begin{aligned}
    R={}&
    \{\overline{v}\mid v\in X\cup Y\cup\{p,n\}\}\\
    &\cup
    \{c_C\mid C\in\mathcal{C}^{+}\}
    \cup
    \{\varphi^{-},t^{+}\}.
\end{aligned}
\]
Jeder Angriff in \(R_{\Psi}\) verläuft zwischen \(L\) und \(R\). Somit gilt
\(
    R_{\Psi}\subseteq(L\times R)\cup(R\times L),
\)
und \(F_{\Psi}\) ist bipartit.

    \begin{figure}[htbp]
        \myfloatalign
        \centering
        \includegraphics[width=0.9\textwidth]{../../Bilder/bipartite_AF-reduction.pdf}
        \caption[Reduktion]{Konstruktionsbeispiel Reduktion}
        \label{fig:bipartite-st-reduction}
    \end{figure}

    \paragraph{Korrektheit.}
    In jedem stabilen Labelling ist aus jedem Paar \(v,\overline{v}\) genau ein Argument
    mit \(\mathsf{in}\) gelabelt. Die Einschränkung eines stabilen Labellings auf die
    Variablenargumente entspricht daher einer totalen Wahrheitsbelegung.

    Angenommen, \(\Psi\) ist gültig. Seien \(\lambda_1\) und \(\lambda_2\) beliebige
    stabile Labellings. Das Labelling \(\lambda_1\) legt die gewünschten Labels von
    \(\varphi^{+}\) und \(\varphi^{-}\) fest, während \(\lambda_2\) eine Belegung von
    \(X\cup\{p,n\}\) festlegt. Da \(\Psi\) gültig ist, existiert zu der durch
    \(\lambda_2\) bestimmten \(X\)-Belegung eine Belegung von \(Y\), welche alle
    ursprünglichen Klauseln von \(\varphi\) erfüllt. Diese Belegung erfüllt die
    Klauseln unabhängig von den Werten der Variablen \(p\) und \(n\). Daher werden
    sämtliche Klauselargumente mit \(\mathsf{out}\) gelabelt.

    Anschließend können die Labels von \(\varphi^{+}\) und \(\varphi^{-}\) unabhängig
    gewählt werden: Für \(\varphi^{+}=\mathsf{in}\) wird \(t^{+}=\mathsf{out}\) gesetzt, während für
    \(\varphi^{+}=\mathsf{out}\) das Argument \(t^{+}\) mit \(\mathsf{in}\) gelabelt wird. Analog wird
    mit \(\varphi^{-}\) und \(t^{-}\) verfahren. Somit existiert ein stabiles Labelling,
    das die Einschränkung von \(\lambda_1\) auf \(A\) und die Einschränkung von
    \(\lambda_2\) auf \(B\) kombiniert.

    Sei nun \(\Psi\) nicht gültig. Dann existiert eine Belegung
    \(\alpha\) von \(X\), sodass keine Belegung von \(Y\) die Formel \(\varphi\)
    unter \(\alpha\) erfüllt. Ergänze diese Belegung durch \(p=\perp\) und \(n=\top\).
    Im zugehörigen stabilen Labelling \(\lambda_1\) sind \(p\) und \(\overline{n}\) mit \(\mathsf{out}\)
    gelabelt.

    Andererseits ist das Labelling
    \(\varphi^{+}=\mathsf{in}\) und \(\varphi^{-}=\mathsf{in}\) als \(\lambda_2\) realisierbar. Hierzu wird
    \(p=\top\) und \(n=\perp\) gewählt. Dann sind \(p\) und \(\overline{n}\) mit
    \(\mathsf{in}\) gelabelt und greifen alle positiven beziehungsweise negativen
    Klauselargumente an. Daher sind sämtliche Klauselargumente \(\mathsf{out}\), sodass
    \(\varphi^{+}\) und \(\varphi^{-}\) gleichzeitig \(\mathsf{in}\) sein können.

    Ein kombinierendes stabiles Labelling müsste gleichzeitig die nicht erfüllende
    Belegung \(\alpha\), die Labels \(p=\mathsf{out},n=\mathsf{in}\) und
    \(\varphi^{+}=\varphi^{-}=\mathsf{in}\) realisieren. Aus
    \(\varphi^{+}=\varphi^{-}=\mathsf{in}\) folgt, dass alle positiven und negativen
    Klauselargumente \(\mathsf{out}\) sind. Da \(p\text{ und } \overline{n}\) keine Klauselargumente
    angreifen, muss jede Klausel von mindestens einem \(\mathsf{in}\)-gelabelten ursprünglichen
    Literalargument angegriffen werden. Die dadurch bestimmte Belegung von \(Y\) würde
    \(\varphi\) unter \(\alpha\) erfüllen, im Widerspruch zur Wahl von \(\alpha\).
\end{proof}

\begin{Proposition}
    \label{prop:bip-prf-ind-hard}
    \(preferred\)-Unabhängigkeit ist auch in bipartiten AFs \(\Pi_2^{\mathrm{P}}\)-vollständig.
\end{Proposition}

\begin{proof}
    Folgt direkt aus Proposition \ref{prop:bip-st-ind-hard} und der Tatsache, dass 
    bipartite Graphen eine Unterklasse der ungeradezykelfreien Graphen sind. Womit direkt
    \(\Lambda_{\mathrm{st}}(F)=\Lambda_{\mathrm{pr}}(F)\) gilt. 
\end{proof}

\iffalse

\section{Symmetrische AFs}
\label{sec:sym-afs}

Symmetrische, irreflexive AFs sind kohärent \textcite{Coste:2005:symmetric-af-irflx-coherent}.

In diesem Abschnitt betrachten wir Argumentationssysteme mit besonders hoher
Angriffsdichte. Die zugrunde liegende Beobachtung ist, dass zusätzliche Angriffe
die Kombinierbarkeit von Labellings einschränken können. Dadurch treten
Verletzungen bedingter Unabhängigkeit häufiger auf. Als Extremfall betrachten
wir zunächst vollständig gerichtete Argumentationssysteme.

\subsection{Vollständig gerichtete Graphen}
\label{subsec:cdg-afs}

Wir nennen ein AF \(F=(A,R)\) \emph{vollständig gerichtet}, wenn für alle
verschiedenen Argumente \(a,b\in A\) sowohl \((a,b)\in R\) als auch
\((b,a)\in R\) gilt. Jedes Argument greift also jedes andere Argument an.

Das folgende Lemma beschreibt die Struktur der \textit{stable}-Labellings in
vollständig gerichteten AFs.

\begin{Lemma}
    \label{lem:cdg-afs-labellings}
    Sei \(F=(A,R)\) ein vollständig gerichtetes AF mit \(A\neq\emptyset\).
    Dann gilt
    \[
        |\Lambda_{\mathrm{st}}(F)|=|A|.
    \]
\end{Lemma}

\begin{proof}
    Für jedes Argument \(a\in A\) definieren wir ein Labelling \(\lambda_a\)
    durch
    \[
        \mathsf{in}(\lambda_a)=\{a\}
        \qquad\text{und}\qquad
        \mathsf{out}(\lambda_a)=A\setminus\{a\}.
    \]
    Da \(a\) jedes Argument aus \(A\setminus\{a\}\) angreift, ist jedes
    Argument mit Label \(\mathsf{out}\) von einem Argument mit Label
    \(\mathsf{in}\) angegriffen. Zugleich wird \(a\) von keinem Argument mit
    Label \(\mathsf{in}\) angegriffen, da \(a\) das einzige Argument mit Label
    \(\mathsf{in}\) ist. Somit ist \(\lambda_a\) ein
    \textit{stable}-Labelling von \(F\). Da jedes \(a\in A\) ein solches
    Labelling liefert, folgt
    \[
        |\Lambda_{\mathrm{st}}(F)|\geq |A|.
    \]

    Umgekehrt sei \(\lambda\in\Lambda_{\mathrm{st}}(F)\). Da
    \(\lambda\) ein \textit{stable}-Labelling ist, gilt
    \(\mathsf{und}(\lambda)=\emptyset\). Außerdem kann
    \(\mathsf{in}(\lambda)\) nicht leer sein: Wäre
    \(\mathsf{in}(\lambda)=\emptyset\), so hätte kein Argument mit Label
    \(\mathsf{out}\) einen Angreifer mit Label \(\mathsf{in}\), im Widerspruch
    zur \textit{stable}-Bedingung.

    Es bleibt zu zeigen, dass höchstens ein Argument mit Label
    \(\mathsf{in}\) existieren kann. Angenommen, es gäbe zwei verschiedene
    Argumente \(a,b\in A\) mit
    \[
        \lambda(a)=\mathsf{in}
        \qquad\text{und}\qquad
        \lambda(b)=\mathsf{in}.
    \]
    Da \(F\) vollständig gerichtet ist, gilt \((a,b)\in R\). Damit wird ein
    Argument mit Label \(\mathsf{in}\) von einem Argument mit Label
    \(\mathsf{in}\) angegriffen. Dies widerspricht der Bedingung, dass ein
    Argument genau dann \(\mathsf{in}\) ist, wenn alle seine Angreifer
    \(\mathsf{out}\) sind. Also enthält jedes \textit{stable}-Labelling genau
    ein Argument mit Label \(\mathsf{in}\).

    Folglich sind die \textit{stable}-Labellings von \(F\) genau die
    Labellings \(\lambda_a\) mit \(a\in A\). Daher gilt
    \[
        |\Lambda_{\mathrm{st}}(F)|=|A|.
    \]
\end{proof}

Insbesondere existiert für jedes \(a\in A\) genau ein
\textit{stable}-Labelling \(\lambda_a\) mit
\[
    \mathsf{in}(\lambda_a)=\{a\}.
\]
Damit kann jedes einzelne Argument in einem \textit{stable}-Labelling
akzeptiert werden, aber niemals gemeinsam mit einem anderen Argument.

\begin{Proposition}
    \label{prop:cdg-afs-not-st-ci}
    Sei \(F=(A,R)\) ein vollständig gerichtetes AF mit \(|A|\geq 2\).
    Seien \(X,Y,Z\subseteq A\) paarweise disjunkt mit
    \(X\neq\emptyset\) und \(Y\neq\emptyset\). Dann gilt
    \[
        X\not\bot_{\mathrm{st}}Y\mid Z.
    \]
\end{Proposition}

\begin{proof}
    Da \(X\) und \(Y\) nichtleer sind, wählen wir Argumente
    \(x\in X\) und \(y\in Y\). Wegen der paarweisen Disjunktheit gilt
    \(x\neq y\) sowie \(x,y\notin Z\).

    Nach Lemma~\ref{lem:cdg-afs-labellings} existieren
    \textit{stable}-Labellings \(\lambda_x,\lambda_y\in
    \Lambda_{\mathrm{st}}(F)\) mit
    \[
        \mathsf{in}(\lambda_x)=\{x\}
        \qquad\text{und}\qquad
        \mathsf{in}(\lambda_y)=\{y\}.
    \]
    Da weder \(x\) noch \(y\) in \(Z\) liegt, erhalten alle Argumente aus
    \(Z\) in beiden Labellings das Label \(\mathsf{out}\). Also gilt
    \[
        \lambda_x|_Z=\lambda_y|_Z.
    \]

    Angenommen, \(X\bot_{\mathrm{st}}Y\mid Z\) würde gelten. Dann müsste es
    nach Definition der bedingten Unabhängigkeit ein
    \(\lambda_3\in\Lambda_{\mathrm{st}}(F)\) geben mit
    \[
        \lambda_3|_X=\lambda_x|_X,
        \qquad
        \lambda_3|_Y=\lambda_y|_Y,
        \qquad
        \lambda_3|_Z=\lambda_x|_Z.
    \]
    Insbesondere folgt daraus
    \[
        \lambda_3(x)=\mathsf{in}
        \qquad\text{und}\qquad
        \lambda_3(y)=\mathsf{in}.
    \]
    Dies ergibt allerdings sofort einen Widerspruch, da \((a,b) \in R\) ist und somit ein mit dem \(\mathsf{in}\)-\emph{Label} versehenes Argument ein weiteres mit
    dem \(\mathsf{in}\)-\emph{Label} versehenes Argument angreifen würde.  

    Also kann kein solches \(\lambda_3\) existieren. Damit gilt
    \[
        X\not\bot_{\mathrm{st}}Y\mid Z.
    \]
\end{proof}

\subsection{Variable Kantenwahrscheinlichkeiten}

\begin{table}[htbp]
    \myfloatalign
    \centering
    \begin{tabular}{l r r} \toprule
        \(p\) & \multicolumn{2}{c}{Singleton-Enumeration} \\
        \cmidrule(lr){2-3}
              & Paare & unabh.\ Anteil \\ \midrule
        \(0{,}1\) & 4\,950 & 90{,}1\,\% \\
        \(0{,}3\) & 4\,950 & 71{,}2\,\% \\
        \(0{,}5\) & 4\,950 & 50{,}5\,\% \\
        \(0{,}8\) & 4\,950 & 19{,}3\,\% \\
        \(0{,}9\) & 4\,950 & 9{,}6\,\% \\
        \(1{,}0\) & 4\,950 & 0\,\%  \\
        \bottomrule
    \end{tabular}
    \caption[Unabh\"angigkeitsrate symmetrischer AFs nach Kantenwahrscheinlichkeit]{Anteil
    unabh\"angiger Anfragen in symmetrischen AFs in Abh\"angigkeit von der
    Kantenwahrscheinlichkeit \(p\). Die Singleton-Enumeration testet
    ersch\"opfend alle \(\binom{n}{2}\) Paare eines Repr\"asentanten (\(n=100\)).}
    \label{tab:gce-sym}
\end{table}
Umso verbundener der Graph, umso geringer der Anteil an unabhängigen Paaren. 


\section{Stark verbundene Komponenten}
\label{sec:scc-afs}

Kaum \emph{stable}-Labellings vorhanden. 


\fi
\fi
